\section{ベイズ推測の基礎2}
\label{D1_lesson26}

\subsection{授業内容}
\subsubsection{科目の中でのこのコマの位置づけ}
\begin{tabular}[t]{cccccc}
    記述統計[4] & 確率[13] & モデル[17] & 推測・推定[8] & 意思決定[6] & 実践[4]
\end{tabular}

\subsubsection{概要}
ベイズの定理によって信念が更新され，確率が変動することを見てきた。今回は連続的な確率，確率分布に対しても同様のロジックで事前分布を事後分布に変えることができることを示す。ここまでの式を改めて連続的な分布関数として考え，それらが統計的推定や，実際の統計モデルでどのように用いられるかを考える。またこうした計算を容易にしたのが確率的プログラミング言語の発展であり，確率分布を乱数によって近似するという方法(MCMC法)の展開である。
\subsubsection{コマ主題細目(箇条書き)}
\begin{description}
\item[事前分布，事後分布]ここまで事前・事後確率と呼んでいたものは，連続的な分布関数出会っても同様に機能する。計算方法として，確率の和が積分になることは，確率質量関数から確率密度関数に変わった時と同じである。また統計モデルにおいてベイズ推定を行う場合，何が更新されるのかイメージが掴みにくいことがある。我々が知りたかったことはパラメータの値であり，これについて割り当てられた確率がモデルによって更新されることを確認する。
\begin{flushright}
    $\to$ \citet{Maeda2017}18-23.や，\citet{豊田201606}のP.18--21
\end{flushright}

\item[尤度] ベイズの定理の分子における，尤度を今一度改めて確認する。特に回帰分析や，要因計画法における尤度関数を考え，どこに事前分布を置き事後分布を得るのかを，改めて確認しておくことが重要である。
\begin{flushright}
    $\to$ \citet{豊田201606}P.17--18や，\citet{kosugi2018}のP.117--122.
\end{flushright}

\item[周辺尤度]ベイズの定理の分母は周辺尤度と呼ばれる。今一度離散変数の例に戻って，周辺尤度の計算方法を考え，翻って連続変数の場合はどのようにして周辺尤度が計算されるかを理解する。この場合は積分の計算が重ねて行われることになるから，近年の技術的なブレイクスルーを見るまではベイズ推測がうまく扱えなかったという経緯がある。
\begin{flushright}
    $\to$ \citet{馬場201907}のP.54
\end{flushright}

\item[MCMC法]マルコフチェイン・モンテカルロ法は，困難であったベイズ推測を可能にした技術的ブレイクスルーであった。基本的なアイデアは，確率分布の形状を求める際に乱数を使って近似するということである。単純な乱数を用いた例から，之による確率分布の近似と積分計算の簡略化ができることを理解し，事後分布の相対的な形状近似ができるようになったことを学ぶ。
\begin{flushright}
    $\to$ \citet{Maeda2017}の7章，\citet{馬場201907}P.62--76
\end{flushright}

\end{description}
\subsubsection{キーワード}
\begin{itemize}
\item 事前分布，事後分布
\item 尤度と周辺尤度
\item MCMC法
\item 乱数による近似

\end{itemize}
\subsection{授業情報}
\paragraph{コマの展開方法}
講義
\paragraph{標準シラバスにおける位置づけ}
科目番号5；心理学統計法；(1)心理学で用いられる統計手法；J より高度な記述や推定を目指して
\subsubsection{予習・復習課題}
\paragraph{予習}
確率を使った推論の理論的基盤であるから，改めて\citet{Maeda2017}のP.16--30を読み，心理統計学，推測統計学を学ぶ理由や心理学における位置づけを確認しておくとよい。ベイズの定理そのものは，簡単な式を展開していくだけでフォローできるが，これをなぜ学ばなければならないのか，明確に学習プロセスに位置付けておくことが重要である。
\paragraph{復習}
ベイズ統計を実践的に使う上での，専門的かつ技術的内容が含まれるのが本時である。MCMCのアルゴリズムなどは計算論として非常に興味深いものであるが，心理学を学ぶ上では「乱数を使って確率分布を近似させることができる」という本質を掴んでおけば十分である。\citet{Maeda2017}は非常にボリュームのある本であるが，5，6，7章を読むとベイズ統計とMCMCの関係が分かり易い。実用的な面については，\citet{豊田201606}のP.29--37がコンパクトにまとまっていて良い。